\documentclass[11pt]{article}
\usepackage[T1]{fontenc}
\usepackage[utf8]{inputenc}
\usepackage{lmodern}
\usepackage[margin=25mm]{geometry}
\usepackage{amsmath,amssymb,amsthm,booktabs,longtable,array,enumitem,microtype}
\usepackage[hyphens]{url}
\usepackage[colorlinks=true,linkcolor=blue!45!black,urlcolor=blue!45!black,citecolor=blue!45!black]{hyperref}
\usepackage{xcolor}
\usepackage{fancyhdr}
\usepackage{listings}
\newcommand{\N}{\mathsf N}
\newcommand{\C}{\mathcal C}
\newcommand{\Con}{\operatorname{Con}}
\newcommand{\code}[1]{\texttt{\detokenize{#1}}}
\newtheorem{proposition}{Proposition}
\newtheorem{lemma}{Lemma}
\setlength{\emergencystretch}{2em}
\setlength{\parskip}{4pt}
\setlength{\parindent}{0pt}
\setlist{itemsep=3pt,topsep=5pt}
\lstset{basicstyle=\ttfamily\small,columns=fullflexible,breaklines=true,frame=single,rulecolor=\color{black!15},xleftmargin=5pt,xrightmargin=5pt}
\pagestyle{fancy}
\fancyhf{}
\fancyhead[L]{\small ProveIt: search below 218 states}
\fancyhead[R]{\small Research checkpoint}
\fancyfoot[C]{\thepage}
\setlength{\headheight}{14pt}
\hypersetup{pdftitle={A search below 218 states for Con(ZFC)},pdfauthor={Research prepared for ProveIt},pdfsubject={Verified 218-state variants, source compression, and scoped negative results}}
\title{A Search Below 218 States for $\Con(\mathrm{ZFC})$\\[5pt]\large Verified variants, alternative calculi, and unresolved reductions}
\author{Research checkpoint prepared for the ProveIt project}
\date{20 October 2026}
\begin{document}
\maketitle
\begin{abstract}
\noindent
No justified machine with 217 or fewer working states was obtained.
This report records the results that survived verification: two new 218-state
variants, an explicit proof-compiler repair for omission of a redundant stack
operation, a 16-instruction Cantor-stack push in place of the existing
19-instruction push, independently checked source invariants, and concrete
obstructions to several attempted state deletions. Two source constructions
reach 204 DAG nodes, but neither supplies the required logical calibration
and valid physical allocation. Large finite searches exclude the documented
families only. The accompanying archive contains literal machine tables,
source files, portable verifiers, certificates, and the experimental sources.
It is a research checkpoint, not a claimed improvement of the 218-state bound.
\end{abstract}
\setcounter{tocdepth}{1}
\tableofcontents
\newpage

\section{Target, baseline, and evidence standard}
The target is a binary, deterministic, single-tape Turing machine $M$ with a
two-way infinite tape, initially all zero, such that
\[
 \mathrm{ZFC}\vdash\bigl(\neg\operatorname{Halt}(M,\mathrm{blank})
 \leftrightarrow\Con(\mathrm{ZFC})\bigr),
\]
using at most 217 working states. Moves are $L$ or $R$. The terminal state
has no outgoing transitions and is not counted. A smaller count in another
machine model would not settle this target.

All experiments start from ProveIt commit
\begin{center}\small\code{9adcaff0a72486ecc81242e0eca86e428c74879e}.\end{center}
The baseline is the repository's \code{candidate-k218}, whose ordinary
mathematical calibration is inherited, not re-proved from the ZFC axioms
in this investigation~\cite{repo}. Its native labels have initial state 0,
dispatcher 84, and terminal label 254. It has 218 working states and 436 rows.
The source has 205 distinct DAG nodes, 4,744 primitive occurrences, maximum
address depth 38, and maximum zero run 6. In the physical allocation, 168
decoder states represent 183 roles, with 15 shared roles; the other 50 states
serve execution and initialization. Thus a one-node source reduction need
not reduce the actual machine by one state.

The archive uses three explicit evidence levels. Exact finite checks establish
literal table identities, finite startup configurations, or counterexamples.
The symbolic runtime engine proves arbitrary-natural unary traversals and
arbitrary finite spare-list traversals by mathematical induction. Neither
kind is a newly accepted Lean or Rocq theorem. Finally, the logical calibration
and source arithmetic arguments are ordinary mathematical proofs; open
calibration obligations are labeled as such. Search failures are never used
as lower bounds.

\begin{center}
\small
\begin{tabular}{p{57mm}rrp{57mm}}
\toprule
Construction & DAG & States & Status\\
\midrule
Pinned baseline &205&218&Inherited calibration; exact baseline files\\
Omit successful-MP dummy push &205&218&Full runtime and blank-start replay; compiler repair below\\
Swap MP and DROP command tags &205&218&Full runtime and blank-start replay; command recoding below\\
Merged NIP/DROP command &204&223&Raw prefix allocation; completeness open\\
K/S-only implication basis &204&---&Calibration open; proposed allocations fail\\
Reversed stack, PUSH16 &206&229&Arithmetic proof; raw allocation only\\
Reassociated source specimen &204&---&Existing decoder marker obligations fail\\
\bottomrule
\end{tabular}
\end{center}
A dash means no valid physical machine was established. The 223- and 229-state
entries are raw allocation counts, not complete machine-correctness claims.

\section{Two fully replayed 218-state alternatives}
\subsection{Native notation and command semantics}
Let
\[
 P(x,y)=\frac{(x+y)(x+y+1)}2+y,
\quad I(A,B)=4P(A,B),
\]
\[
 R(a,b)=4P(a,b)+1,\quad Q_aB=4P(a,B)+2,\quad E(a,b)=4P(a,b)+3.
\]
Write $F=R(0,0)=1$ and $\N A=I(A,F)$. The designated theorem code 0 is
interpreted as truth. These conventions are native to the inherited calculus;
raw natural codes and numeric binder labels require their own calibration.
The work tuple is $(t,S,d,h)$: current theorem, theorem stack, additive
parameter, and cached theorem. The stack is Cantor-coded, with code 0 behaving
as an infinite tail of theorem 0.

For $B=d+x$, successful MP recognizes $h=I(t,B)$. Its baseline result is
$(B,P(0,S),B,0)$. A failed MP returns $(0,S,B,0)$.
DROP(0) maps $(t,S,d,h)$ to $(h,P(t,S),d,h)$.
NIP(0) halts if $h=1$; otherwise it pops the next stack theorem into $h$,
leaving $t$ unchanged. AX emits nine proved formulas, followed on the stack
by the previous current theorem and previous stack. Its current result is a
gated set-theoretic bundle or 0. The command parameter order is
\code{Cmd.ax c x y u}; in particular, \code{ax B 0 0 0} sets $c=B$,
not $x=B$.

\subsection{Omitting the dummy zero after successful MP}
Change only the successful branch to
\[
 (t,S,d,h)\longmapsto(B,S,B,0),\qquad B=d+x.
\]
The source edit deletes the 19-leaf PUSH subtree at address \code{111110}
and promotes its seven-leaf COPY sibling at \code{111111} to
\code{11111}. The packaged verifier reconstructs these exact bytes.
There are now 4,725 primitive occurrences but still 205 DAG nodes.
Only native transition row $164/1$ changes.

\begin{proposition}
Under the inherited source semantics and cofinal command enumeration, the
modified enumerator halts exactly when the same native calculus derives $F$.
\end{proposition}
\begin{proof}
Soundness follows from the same invariant: all current, cached, and stacked
formulas are theorems. A successful MP now leaves a sound stack unchanged.
For completeness, use the existing compiler contract: from arbitrary current
value $T$ and stack $S$, with $d=h=0$, the compiled derivation either halts or
returns $(X,S,0,0)$ for its conclusion $X$.

The leaf and generalization blocks are unchanged. The discard block forces
MP failure: it compares the zero cache with $I(T,1)>0$. Thus it never depends
on the modified branch. Only the MP compiler's cleanup changes. After the
successful new MP, its state is $(B,S,B,0)$. The original tail was
\[
 \operatorname{AX}(B,0,0,0);\ \operatorname{NIP}^{10};\
 \operatorname{DROP}(0);\ \operatorname{NIP}(0);\ \operatorname{NIP}(0).
\]
Delete its final NIP. If $L$ denotes the nine emitted theorem entries,
the new tail has the exact trace
\begin{align*}
 (B,S,B,0)&\longmapsto(0,L::B::S,0,0)\\
 &\longmapsto(0,S,0,B)\\
 &\longmapsto(B,0::S,0,B)\\
 &\longmapsto(B,S,0,0).
\end{align*}
If a NIP sees $F$ earlier, it already halts soundly. In particular, $B=F$
halts at the final step instead of returning, as the disjunctive compiler
contract permits. Induction on the derivation therefore transfers with this
single cleanup deletion. Finite program packing and cofinal enumeration
are unchanged. The reference definitions are \code{Contract},
\code{discard}, \code{mpBlock}, and \code{mp_contract} in
\code{Enum222.lean}, included for reference.
\end{proof}
This is a complete local repair of the inherited ordinary compiler argument;
no modified Lean theorem is claimed. In particular, the source is intentionally
\emph{not} ordinal-trace-equivalent to the old source. An early receipt's
inherited wording suggesting that obligation was corrected before packaging.

\subsection{Swapping MP and DROP tags}
The second variant retains every command body and uses tag 0 for DROP, tag 1
for MP, and tag 2 for NIP; tags at least 3 retain the axiom command. The baseline
uses MP, DROP, NIP in that order. The closed source bodies at addresses
\code{111}, \code{110101}, and \code{110100101} are permuted literally.
Their selector scratch register is zero on entry. The verifier checks the
closed-boundary conditions, both control-flow constructions, and the exact
source bytes. Only native rows $135/1$ and $185/1$ change.

This gives a computable bijection on finite command programs, preserving
command parameters and command semantics. Every compiled refutation therefore
has a recoded finite program. The inherited reset/padding construction gives
arbitrarily large encodings of such a refutation: vary the fourth parameter
of a prepended AX, whose output consists entirely of theorems and which resets
$d=h=0$. That parameter participates in the emitted classical axiom, but the
subsequent compiler accepts the resulting arbitrary current theorem and stack.
Thus the enumeration tail from either initial counter eventually examines a
refuting program. Soundness is unchanged. The blank computation reaches a clean seed with
candidate counter 2 rather than 3; the inherited safe-start/cofinality argument
covers either seed. Thus initialization time changes without changing the
finite-refutation criterion.

\subsection{Exact operational certificates}
The bundled unchanged runtime engine was replayed with a fresh portable driver
against both variants. It proves every primitive interface for independent
arbitrary natural primary registers, extra gap, first spare value, and an
arbitrary finite remaining spare list. The endpoint comparison includes the
entire tape, head, and state. The bound 4,000 limits proof construction, not
the represented values or concrete execution time. An independent audit
checks complete ordered guard coverage, positive affine durations, all literal
row coverage, and rederives the four unary-block schemas directly from rows.

\begin{center}
\begin{tabular}{lrrr}
\toprule
Variant & Guarded interfaces & Symbolic rules & Minimum duration\\
\midrule
Omit dummy push &6,142&1,511,359&12\\
Swap MP/DROP (DMN) &6,167&1,516,795&12\\
\bottomrule
\end{tabular}
\end{center}
Both use all 436 rows. A second implementation, using a dense two-sided tape,
independently verifies the complete blank endpoint against a directly
constructed frame, without invoking the original frame decoder:
\begin{center}
\begin{tabular}{lrrrrr}
\toprule
Variant & Steps & Head & Extra gap & Unit spares & One cells\\
\midrule
Baseline &312,909&$-5$&80&130&277\\
Omit dummy push &312,909&$-5$&80&130&277\\
DMN &123,776&$-4$&23&72&159\\
\bottomrule
\end{tabular}
\end{center}
Every endpoint is native dispatcher 84 at root source address. Baseline and
omission have logical vector $(3,0,0,0,3,0,0,0,0)$; DMN has
$(2,0,0,0,2,0,0,0,0)$. All cells outside the recorded finite support are zero.
Canonical tables relabel working states to $0,\ldots,217$ and halt to 218;
the archive contains the exact bijections.

\input{sections/formalism}
\input{sections/source}
\input{sections/runtime}

\section{What would actually close the remaining gap?}
The computational evidence favors changes to the shared source and its physical
representation together. Reducing primitive occurrences alone is insufficient:
PUSH16 saves 279 occurrences but loses enough decoder sharing to increase the
raw state count. A 204-node DAG is also insufficient: the merged command
requires a new completeness compiler and the K/S-only proposal requires a
calibration theorem, while both need a valid physical construction.

Three concrete unresolved tasks remain worth distinguishing.
\begin{enumerate}
\item \textbf{Calibrate the actual K/S-only generated calculus.} A proof using
all equality rows, commuting quantifiers, and the negative set kernel could
settle whether its 204-node source has the needed consistency strength.
The small Heyting countermodel rules out only the proposed modal induction;
it does not model the complete set calculus. A successful argument would
still need a new physical allocation.
\item \textbf{Compile proofs for the merged command.} A continuation-based
compiler that preserves the desired premise while crossing AX's nine extra
outputs would address its exact completeness obstruction. The current
223-state allocation would then need six further physical savings.
\item \textbf{Redesign the decoder's stored context.} The original 218-state
improvement reused a preceding tape marker. The current negative search
families restrict local joins, offsets, and helper changes. A different
invariant or representation could permit a reduction outside those families.
The provided ternary certificate gives a verified starting domain, but it
is an overapproximation of concrete reachability.
\end{enumerate}
Other calculi offer conceptual alternatives, but their specification size is
not their binary-machine cost. Pfenning's analysis of short implicational
bases informed the basis substitutions~\cite{pfenning}. Andr\'eka and
N\'emeti give a translation of set theory into weak logics with three
commuting modalities~\cite{andreka}; no competitive source encoder for that
translation was constructed here. Neither reference supplies an unmeasured
state-count saving.

\section{Archive map and reproduction}
The archive is self-contained for its stated new checks. It does not contain
a complete Lean build of the inherited repository, and copied reference files
are not advertised as such a build.
\begin{center}
\small
\begin{tabular}{p{67mm}p{84mm}}
\toprule
Path & Contents\\
\midrule
\code{variants/omit-zero/} & Native and canonical 218-state tables, exact source, original receipts\\
\code{variants/dmn/} & Native and canonical 218-state tables, command-permuted source, original receipts\\
\code{code/} & Portable fresh runtime replay and independent receipt/source audit\\
\code{vendor/portable/} & Unchanged, hash-pinned parent runtime and source utilities\\
\code{evidence/} & Fresh receipts produced by the packaged driver and audit\\
\code{experiments/formalism/} & Formula experiments, finite countermodels, exact compiler notes\\
\code{experiments/source_compression/} & PUSH16, 204-node sources, prefix-code sweep, invariant certificates\\
\code{experiments/machine_compression/} & Rejected 217-state tables, concrete witnesses, search scopes and source\\
\code{baseline/}, \code{references/} & Pinned inherited table/source and relevant repository definitions\\
\code{MANIFEST.json} & SHA-256 for every packaged non-manifest file\\
\bottomrule
\end{tabular}
\end{center}
Paths in the detailed experimental sections are relative to the corresponding
checkpoint directory. Each checkpoint has its own README. The root manifest
is authoritative for the complete package.

From the extracted archive root, run:
\begin{lstlisting}
python3 code/verify.py
python3 code/replay_runtime.py omit-zero
python3 code/replay_runtime.py dmn
python3 experiments/source_compression/code/verify_checkpoint.py
python3 experiments/source_compression/code/verify_invariants.py
\end{lstlisting}
The first command audits the saved receipts and reruns exact finite checks.
The next two rederive all symbolic interfaces. The source commands reconstruct
selected edits and check the invariant certificates; the ternary check needs
a C++17 compiler. Use Python 3 with assertions enabled. These commands need
no Python packages or network access. Optional searches have their own limits.

Rebuild this report by running \code{pdflatex} twice on
\code{zfc-217-state-search.tex}; all TeX fragments are included.

\begin{thebibliography}{9}
\small\raggedright
\setlength{\itemsep}{2pt}
\setlength{\parskip}{0pt}
\bibitem{repo}
Vladimir Reshetnikov and ProveIt contributors.
\emph{ProveIt: candidate-k218}, pinned commit
\code{9adcaff0a72486ecc81242e0eca86e428c74879e}, 2026.
\url{https://github.com/VladimirReshetnikov/ProveIt/tree/9adcaff0a72486ecc81242e0eca86e428c74879e/Computability/BusyBeaver/Research/ZFC/candidate-k218}.
See the included upstream proof summary and provenance.
\bibitem{pfenning}
Frank Pfenning. \emph{Single Axioms in the Implicational Propositional Calculus}.
CADE, 1988. \url{https://www.cs.cmu.edu/~fp/papers/cade88.pdf}.
\bibitem{andreka}
Hajnal Andr\'eka and Istv\'an N\'emeti.
\emph{Formalizing set theory in weak logics, searching for the weakest logic with G\"odel's incompleteness property}.
arXiv:1111.0995, 2011. \url{https://arxiv.org/abs/1111.0995}.
\end{thebibliography}
\end{document}
